\section{Artin底上の反標準section再構成}
\label{sec:cubic-formal-intrinsic-reconstruction}

\begin{theorem}\label{thm:h3-artin-intrinsic-section-reconstruction}
$(\mathcal A,\mathcal O)$ を幾何学的H3対象とし、$s=S^{-1}[2]$ とする。
剰余体 $k$ の可換局所Artin環 $A$ 上で $\mathcal A$ の
Lowen--Van den Berghの意味での平坦なAbel圏変形 $\mathcal A_A$ を取り、
$s$ がexact自己同値 $s_A$ に持ち上がると仮定する。
このとき $\mathcal O$ には平坦な持ち上げ $\mathcal O_A$ が同型を除き一意に存在し、
\[
 R_A=\bigoplus_{n\ge0}
 \operatorname{Hom}_{\mathcal A_A}(\mathcal O_A,s_A^n\mathcal O_A)
\]
は次数ごとに有限自由で、Artin底変更と可換する。
三つ組 $(\mathcal A_A,\mathcal O_A,s_A)$ は
$(\qgr R_A,\pi R_A,(1))$ と同値である。
$R_A$ は一次生成で、その関係は次数2と3で生成される。
逆に、反標準環の全次数平坦な次数付きArtin変形は、
この形の構造対象とliftしたshiftを持つ平坦なAbel圏変形を与える。
特に、相対Serre作用の逆から $s_A$ が得られる場合、
section環の全次数平坦性と豊富性を別に仮定する必要はない。
\end{theorem}
\begin{proof}
幾何学的H3では $\mathcal O$ がexceptionalなので、
\cite[命題5.1]{VdBFormalAbelian}の持ち上げ判定が成立する。
$s_A^n\mathcal O_A$ も平坦である。
幾何学的fiberで $\operatorname{Ext}^1(\mathcal O,s^n\mathcal O)=0$
は正規三次商のAS Cohen--Macaulay性による。
従って \cite[命題5.3]{VdBFormalAbelian}を
各二つの平坦対象に適用して、sectionの自由性・底変更を得る。
負の次数のsectionは同じ議論により零、次数0は単位写像で $A$ と同定される。

幾何学的fiberの $\mathcal O(n)$ は同文献の意味でstrongly ampleである。
生成はH3の豊富性から従う。
任意の $F$ への $\mathcal O(-m_i)$ の有限直和からの全射を取り、
Serre双対を使うと
$\operatorname{Ext}^2(\mathcal O(-n),F)^*=\operatorname{Hom}(F,\mathcal O(-n-1))=0$
が大きい $n$ で成り立つ。実際、右辺は全射から負次数のsectionの直和に入る。
全射の核に同じ二次消滅を適用すると、長完全列から
十分大きい $n$ で $\operatorname{Ext}^1(\mathcal O(-n),F)=0$ も従う。
三次以上は幾何学的fiberの大域次元2で消える。
従って \cite[定理4.2]{VdBFormalAbelian}で
$\mathcal O_A(n)$ のstrong amplenessと圏のExt有限性が従い、
Artin--Zhang再構成を適用できる。$s_A$ があるため、
得られる周期的な$\mathbb Z$代数は上の次数付き環から来る。

一次生成は各次数の乗法写像を分裂する自由加群上で比較し、
剰余体上の全射からNakayama補題で従う。
$J_n$ を、$R_{A,1}^{\otimes n}\to R_{A,n}$ の核とする。
右辺が自由なので核の形成も剰余体化と可換する。
剰余体上の全関係は次数2,3で生成されるから、
$J_n$ を $J_2,J_3$ の張る二側イデアルの次数 $n$ 部分で割った有限加群は
剰余体上で零、従って零である。
こうして任意のArtin平坦圏変形でsection環の低次数提示が得られる。

逆に、全次数自由な次数付き変形 $R_A$ を与える。
$A$ の極大イデアルの冪による随伴環は $R\otimes_k\operatorname{gr} A$ であり、
有限長のフィルターから $R_A$ はnoetherianである。
次数付き加群圏は、次数付き自由shiftの小線型圏の加群圏である。
そのHomは全て $A$ 上自由なので、
\cite[定理8.16]{LowenVdBFlatAbelian}の平坦な加群圏変形となる。
torsion商はその局所化であり、
$A$ の極大イデアルに消される加群上では元のtorsion商に一致する。
実際、$R_{A,\ge n}$ からそのような加群への射は
剰余体にした $R_{\ge n}$ からの射に等しいので、飽和の式が一致する。
局所化の右随伴は入射対象を入射対象へ送る。
有限 $A$ 加群からの内部 $\operatorname{Hom}_A$ は有限表示の核で計算でき、
exactな商関手と可換するので、商圏の入射対象もcoflatである。
従って商圏は平坦であり、noetherian部分を取ることは
\cite[定理7.1, 8.8]{LowenVdBFlatAbelian}と両立する。
自由加群の像 $\pi R_A$ は底上平坦で、次数移動がliftである。
上のsection再構成を適用すると、sectionへの自然写像は自由加群間の写像で
剰余体上同型だから全次数で同型である。
\end{proof}

\begin{corollary}\label{cor:h3-artin-finite-presentation-ring-hull}
全次数平坦な二次包絡も持つものに限ると、
上のsection変形の三次超曲面表示は各幾何学的H3点で
有限型incidenceの形式完備化によって表される。
ここで一次基底を固定する。基底を除くと表示の群oidは基底変更商であり、
shiftに可換する構造対象付き矢だけを取る場合は共通次数スカラーを除く。
\end{corollary}
\begin{proof}
$V=A^4$ とし、二次の全関係 $J_2$、三次の全関係 $J_3$ を持つ
前掲の有限型Grassmann incidenceを取る。
$B=T(V)/(J_2)$、$R=T(V)/(J_2,J_3)$ について、
各次数の関係挿入行列にそれぞれ所望階数の上界を課す。
全次数の小行列式を使っても、noetherian性から閉条件のイデアルは有限生成である。
幾何学的H3点を固定したArtin局所変形では、各次数の階数下界は
剰余体で非零のminorが単元になることから自動である。
従って両環の全次数平坦性とこの閉条件は形式近傍で同値となる。

$B_3\to R_3$ の核は階数1の束で、剰余体の切断へ写る。
その生成元 $G$ の左右乗法は各次数で剰余体上単射だから、
自由加群の分裂単射である。
剰余体で全核を生成するので、Nakayama補題と両辺の平坦性から
全次数で左右とも全核を生成する。
従って形式表示も正則正規三次商である。
相対Koszul性は、有限の内部次数ごとに、augmentationのbar複体を調べれば従う。
その各項は自由で、剰余体上でホモロジーがKoszul対角の一箇所に集中する。
単元minorで分裂することにより、同じ集中と自由性がArtin底上で成立する。

構造対象の全整数shift間の射はsection環から復元され、
これらは有限表示で全対象を生成する。
従って前定理の同値はshiftに可換する矢にも適用できる。
比較同型の変更は $\operatorname{End}(\mathcal O_A)=A$ と、
隣接するshift間の $R_{A,1}$ の忠実性から共通次数スカラーとなる。
同型方程式は一次の基底変更が $J_2,J_3$ を保つ条件であり、有限である。
\end{proof}

全次数平坦な二次包絡を持つことが、任意の平坦な圏変形から
自動的に従うとはここでは主張しない。
また相対Serre作用の存在とそのheartへの制限は、一般の底上では別の認識問題である。

\subsection{相対局所双対のorientationと中山作用}
\begin{proposition}\label{prop:h3-artin-relative-dualizing-line}
系\ref{cor:h3-artin-finite-presentation-ring-hull}の相対表示を取る。
$W_4$ を包絡のKoszul top直線、$\mathcal G$ を三次切断直線とし、
$\mathcal L_{\rm AS}=W_4\otimes\mathcal G^\vee$ と置く。
$Gx=\sigma(x)G$ で正常化を定め、
$\mu=(\mu_B\sigma)|_R$ とする。
正次数イデアルに関する右局所コホモロジーは次数3に集中し、その相対次数付き双対は
\[
 \bigl(H^3_{R_{\ge1}}(R)\bigr)^\vee
 \ \simeq\ {}^\mu R^1(-1)\otimes_A\mathcal L_{\rm AS}
\]
である。ここで $\vee$ は各内部次数の $A$ 双対である。
自然な積対合とこの同型はArtin底変更と可換する。
\end{proposition}
\begin{proof}
相対Koszul双対のtopへの積は完全対合であり、
そのFrobenius中山作用は相対的に定義される。
これを用いたKoszul双加群解消の双対は、長さを反転した同じ解消の
内部次数4のshiftと中山twistである。
この計算は、$W_i$ と $W_{4-i}$ の対合
$W_i\otimes_A W_{4-i}\leftarrow W_4$ の双対
（Koszul双対側では積）で二つの端の微分を比較するもので、
有限自由の各項で体上のKoszul双対計算と同じ式になる。
従って右加群 $M$ が内部次数を合わせて $A$ 上有限自由なら、
相対双加群解消を $M$ とテンソルしてから双対を取ると
\[
 \mathbf R\underline{\operatorname{Hom}}_B(M,B)
 \simeq M^\vee(4)_{\mu_B}\otimes_A W_4^\vee[-4]
\]
を得る。項ごとの同定は
$\operatorname{Hom}_A(M\otimes_A W_i,B)
 \simeq B\otimes_A W_i^\vee\otimes_A M^\vee$ である。
$M$ に左作用もある場合、式の右作用はその前合成から来る。

$M=R/R_{\ge n}$ とする。右 $B$ 加群としての正則三次解消は
$\mathcal G\otimes_A{}^{\sigma^{-1}}\!B(-3)\xrightarrow{G\cdot}B$ であり、
$\mathbf R\operatorname{Hom}_B(R,B)$ は
$R(3)_{\sigma^{-1}}\otimes_A\mathcal G^\vee[-1]$ となる。
変換の随伴と、$M$ 自身の $\sigma$ による半線型同型を使うと、
\[
 \mathbf R\underline{\operatorname{Hom}}_R(M,R)
 \simeq M^\vee(1)_{\mu_B\sigma}\otimes_A\mathcal L_{\rm AS}^{\vee}[-3]
\]
を得る。右作用を確かめるには、$\sigma^{-1}$ 付き標的への射 $f$ を
$f\circ\sigma$ に移すと、その前合成の右作用が $\sigma^{-1}$ に捩れることを見ればよい。
従って変換後の元の右作用には $\sigma$ が掛かる。
この式は $R/R_{\ge n}$ の間のすべての標準射に対し自然である。
局所torsionの標準直極限を取ると、各内部次数で有限自由成分が安定し、
$R^\vee(1)_{\mu}\otimes_A\mathcal L_{\rm AS}^\vee$ を得る。
その双対は主張の式である。
\end{proof}

このorientation直線は、共通次数スカラーに関して重み $4-3=1$ を持つ。
従って命題\ref{prop:h3-presentation-canonical-degree-descent}の
中立化では、$\det V\otimes\mathcal G^\vee$ の代わりに
$\mathcal L_{\rm AS}$ を使うこともできる。

\begin{corollary}\label{cor:h3-artin-relative-nakayama-scalar}
このArtin表示が内在圏の相対Serre作用により
$(1)=S_A^{-1}[2]$ と再構成されるとする。
すると $\mu=(\mu_B\sigma)|_R$ は共通次数スカラーである。
\end{corollary}
\begin{proof}
命題の局所双対から、$\mathcal O_i=\pi R(i)$ 間の第二Extは
$H^3_{R_{\ge1}}(R)$ の対応する内部次数であり、
次数零Homはsection環の対応次数である。
特に $\operatorname{Ext}^2(\mathcal O_i,\mathcal O_{i-1})$
は $\mathcal L_{\rm AS}^\vee$ に正則に同定される。
この直線を自明化する。相対Serreのtraceは、各 $\mathcal O_i$ 上で
この直線を $A$ に移す単元倍の写像である。
実際、$\operatorname{End}(\mathcal O_i)=A$ と完全対合をなす。

$b\in\operatorname{Hom}(\mathcal O_i,\mathcal O_{i+1})=R_1$ と、
$R_1^\vee$ に対応する
$\varphi\in\operatorname{Ext}^2(\mathcal O_{i+1},\mathcal O_{i-1})$ を取る。
二つのYoneda積の値をtraceで比較する。
局所双対の双加群作用により、一方は $\varphi(b)$、
もう一方は $\varphi(\mu(b))$ となり、Serreの自然性により
それぞれのtraceの単元倍を取ると一致する。
従って $\mu|_{R_1}$ はその二単元の比によるスカラー倍である。
一次生成性により全次数の作用は共通次数スカラーである。
\end{proof}

\begin{theorem}\label{thm:h3-artin-relative-nakayama-trivial}
前系の仮定の下で、非被約Artin底上でも
\[
 \mu=1,\qquad \sigma=\mu_B^{-1}
\]
である。従って反標準性が課す正規三次条件は
Artin底上でも単一のNakayama線型核で与えられる。
\end{theorem}
\begin{proof}
$X=K_B\otimes_B R$ とする。ここで $K_B$ は右augmentation加群 $A$ の
有限自由Koszul解消である。
$X$ は有限自由な右 $R$ 複体で、コホモロジーは
$A$ とその三次shiftの二項のみなので、次数torsionである。
特に $X$ の内部次数0の複体は $A$ 一項である。
次数0の導来自己準同型に対し、その内部次数0への制限は
恒等射で与えられる $A\to\mathbf R\operatorname{Hom}_{\mathrm{Gr}R}(X,X)$ を分裂させる。
この分裂は有限自由複体上で直接定義される。

局所双対をこのtorsion完全複体に適用する。
$D_A$ を $A$ 上の導来双対とすると、任意の次数付き完全複体 $Y$ に対し
\[
 D_A\,\mathbf R\operatorname{Hom}_{\mathrm{Gr}R}(X,Y)
 \simeq
 \mathbf R\operatorname{Hom}_{\mathrm{Gr}R}
   \bigl(Y,X\otimes_R({}^{\mu}R^1(-1)\otimes_A\mathcal L_{\rm AS})[3]\bigr).
\]
実際、有限自由な $X$ を用いたtensor--Homの公式で、
左辺は $D_A\,\mathbf R\underline{\operatorname{Hom}}_R(X,R)$ を標的とする
右Homに移る。$X$ はtorsionなので、
$\mathbf R\underline{\operatorname{Hom}}_R(X,R)$ では
標的を局所torsion複体に置き換えられる。
これはtorsion商への導来随伴で、商における $X$ が零であることから従う。
有限自由項で双対化し、命題\ref{prop:h3-artin-relative-dualizing-line}を使うと
右辺の標的となる。以上の同定には内部次数付きの複体を使うので、
内部shiftとの比較は自明である。

$\mathcal L_{\rm AS}$ を自明化し、$\mu=\xi_c$ と書く。
上の双対式のSerre標的は $F=[3](-1)\Phi$、$\Phi=T_{\xi_{c^{-1}}}$ である。
トレースを恒等射の双対として取ると、上記の分裂により
$\operatorname{Hom}_{\mathrm{Gr}R}(X,FX)\to A$ は全射である。
通常の graded trace の計算
（\cite[Lemma 2.6, Theorem 4.1の証明]{RRZtri}）をここで使う。
この計算に必要なのはトレースの自然性と、複体のshiftでのKoszul符号だけである。
本件では有限自由複体の $A$ 双対で同じ恒等式を得るので、
全Ext加群が自由であることを仮定しない。
具体的には、$\Phi$ を恒等と比較する写像は
内部次数 $n$ の元に $c^{-n}$ を掛ける。
従って整数 $j$ の内部shift後の比較は、shift前の比較に比べ $c^j$ 倍である。
具体的に $U=[3](-1)$、$u_X:UX\to FX$ をこの比較とし、
$a:X\to UX$ を取る。
トレースの自然性と比較の自然性により
\[
 \operatorname{Tr}_X(u_Xa)
 =\operatorname{Tr}_{UX}(F(a)u_X)
 =\operatorname{Tr}_{UX}(u_{UX}U(a)).
\]
内部shiftは $-1$ なので $u_{UX}=c^{-1}U(u_X)$ であり、
コホモロジーshift3の符号は $(-1)^{3(3+1)}=1$ である。
従って最後の項は $c^{-1}\operatorname{Tr}_X(u_Xa)$ となる。
トレースが $1\in A$ を取るので $c=1$ である。

$R_1=B_1$ であるから $\mu=(\mu_B\sigma)|_R=1$ は
$\mu_B\sigma$ が包絡の一次空間上恒等であることを意味する。
一次生成性で包絡全体でも恒等となり、線型条件を得る。
\end{proof}

\begin{theorem}\label{thm:h3-artin-relative-serre-exists}
系\ref{cor:h3-artin-finite-presentation-ring-hull}のArtin表示について、
$\mathcal P_A$ を $D(\operatorname{QGr}R)$ のうち
$\pi R(i)$ 全体で生成される厚い部分圏とする。
全ての $A$-平坦なcoherent対象は $\mathcal P_A$ に入る。
この圏の $A$ 相対Serre関手は
\[
 S_A=T_{\mu^{-1}}(-1)\otimes_A\mathcal L_{\rm AS}[2]
\]
である。従ってその逆の二次shiftは $A$-平坦なheart対象を保つ。
正規三次がNakayama線型条件を満たすなら、
orientation直線を除いた反標準作用は表示のshiftになる。
\end{theorem}
\begin{proof}
まず $A$-平坦なcoherent対象 $F$ を取る。
$R$ の商から二次包絡へのexact忘却を $i_*$ とする。
$i_*F$ も底上平坦である。
幾何学的AS4包絡のshift列はstrongly ampleで、
Artin底上も \cite[定理4.2]{VdBFormalAbelian} により同じである。
従って十分高い次数のsectionで $i_*F$ を表す有限生成次数付き加群 $N$ を作れ、
その各成分は同文献の命題5.3により $A$ 上自由である。
剰余体にした $N_k$ の最小次数付き自由解消を持ち上げる。
各段の核は底上自由で、次の自由生成元は剰余体から持ち上がる。
幾何学的包絡の大域次元は4なので有限段で核が自由となり、
剰余体の基底を持ち上げて元の核との同型を得る。
従って $i_*F$ は有限自由shiftの像で解消される。

一方、$F$ の底上平坦性と極大イデアルの有限濾過により、
任意のcoherent $H$ と $j>2$ に対し $\operatorname{Ext}^{j}_{\qgr R}(F,H)=0$ である。
各濾過商が剰余体で消される場合のExtは
平坦な第一引数による底変更で剰余体上のExtに等しく、
残りは有限個の長完全列による。
従って正規三次の導来自己交叉の変換射の二乗は零である。
定理\ref{thm:eta-square-H3}と同じ二つの錐を取ると、$F$ は
二次包絡の完全複体を制限して得る複体のextensionの直和因子である。
従って $F\in\mathcal P_A$ である。

次に右局所torsion複体を $L=\mathbf R\Gamma_{R_{\ge1}}R$ と書く。
局所コホモロジーの次元は有限であることを確認する。
包絡の $B/B_{\ge n}$ は $A$ 上有限自由なので、
包絡の双加群Koszul解消から射影次元は高々4である。
商 $R$ への変換で $\operatorname{Tor}_1^B(B/B_{\ge n},R)$ が現れるが、
正則三次の二項解消により、内部次数は $n-3$ 以上にしかない。
従って $n$ を三つより多く上げた標準射ではこのTorへの写像は零となる。
変換の三角に $\mathbf R\operatorname{Hom}_R(-,M)$ を適用し直極限を取れば、
$R$ 加群 $M$ の局所コホモロジーは包絡側で計算でき、次元は高々4である。
左でも同じである。

各 $B/B_{\ge n}$、$R/R_{\ge n}$ は有限提示で、各次数での射影解消も
有限生成項に取れるため、導来局所torsionは直和と可換する。
したがって右加群の導来局所torsionは $-\otimes_R^{\mathbf L}L$ で与えられる。
この比較は次数付き自由加群で同型で、一般の加群については自由解消を取り、
固定コホモロジー次数で有限コホモロジー次元の範囲まで比較すればよい。
左にも同じ比較がある。
次数torsion加群の次数付き入射包もtorsionである。
非torsion斉次元が入射包にあれば、その巡回部分と元のtorsion加群との交叉は
noetherian性から次数の上に有界となり、高次部分が交叉しないので本質拡大性に反する。
従って導来局所torsionはtorsion加群上で恒等である。
左右の局所torsion複体は、順に左右torsionを取った
二重tensorで比較できる。両者のコホモロジーは
命題\ref{prop:h3-artin-relative-dualizing-line}とその左版により両側torsionなので、
この比較から左右は一致する。

$D$ を三角 $L\to R\to D\to L[1]$ の局所化複体とし、
$J=D_A L={}^{\mu}R^1(-1)\otimes_A\mathcal L_{\rm AS}[3]$ と書く。
局所化は $-\otimes_R^{\mathbf L}D$ であり、$D\otimes_R^{\mathbf L}D\simeq D$。
双対した三角の射 $R^\vee\to J$ は、$R^\vee$ がtorsionであるため
$\mathbf R\Gamma_{R_{\ge1}}J$ を経由する。
この後者は $J$ の可逆双加群部分を局所torsionに掛けて $R^\vee$ に同定される。
経由する自己準同型は同型である。
実際、右augmentation加群 $A$ からのHomを取ると、この射は
局所torsion等式 $A\otimes_R^{\mathbf L}L\simeq A$ の双対である。
従ってsocleの $A$ 上で単元である。
次数付き $R^\vee$ の次数0自己準同型は、
各次数の有限自由双対で左正則加群 $R$ の次数0自己準同型に対応するため、
この単元判定で全体の同型が従う。
よって双対三角は $J$ の局所化三角と同定され、
\[
 D_A D\simeq ({}^\mu R^1(-1)\otimes_A\mathcal L_{\rm AS})[2]
                     \otimes_R^{\mathbf L}D
\]
が得られる。

最後の式は局所化されたshift対象のdg Homにも適用できる。
実際、自由shiftの次数付きdg範疇からそのtails像への射は、
左右に $D$ を掛けると冪等となる。
従ってtailsの双加群の制限は完全忠実であり、
上の両側局所な双加群同型はtailsのdg双加群同型を与える。
各内部次数のHom複体は $H^0,H^2$ が有限自由なので $A$ 上完全である。
このdg双加群の $A$ 双対を用いたtensor--Hom対合は、
有限coneと直和因子で生成される $\mathcal P_A$ 上の相対Serre対合となる。
左中山twistを右加群へのtensorに移すと逆twist、内部shiftは $-1$、
コホモロジーshiftは2となり、主張の式を得る。
\end{proof}

\begin{corollary}\label{cor:h3-artin-intrinsic-formal-atlas}
全次数平坦な二次包絡を持つArtin変形について、
相対Serre作用から定まる反標準section表示の関手は、
有限型の正規三次incidenceの形式完備化によって表される。
一次基底を忘れた構造対象付き同値の群oidでは $\mathrm{PGL}_4$ で割る。
\end{corollary}
\begin{proof}
系\ref{cor:h3-artin-finite-presentation-ring-hull}の両Hilbert上界の
有限型incidenceを取る。Koszul双対の次数3,4の有限階数stratum上で、
Frobenius対合の行列から $\mu_B|_{B_1}$ が底変更と両立して定まる。
三次の核直線に対し、$B_1$ に関する四つの式
\[
 Gx=\mu_B^{-1}(x)G
\]
を課す。これは $B_4$ の有限自由部分上の有限個の閉方程式である。
任意のArtin変形の形式近傍では両方のHilbert下界および
必要なKoszul双対の階数下界は単元minorで保証される。
従ってこの有限型incidenceの形式完備化で、全次数平坦な表示とこの方程式が表される。

相対Serre付き内在側からは、定理
\ref{thm:h3-artin-intrinsic-section-reconstruction}で圏とsectionを復元し、
定理\ref{thm:h3-artin-relative-nakayama-trivial}で閉方程式を得る。
逆にこの方程式を満たす表示では、定理
\ref{thm:h3-artin-relative-serre-exists}で相対Serre双対が存在し、
反標準作用は、底のorientation直線を除けば次数移動である。
Artin局所底ではこの直線は自明であり、自明化の変更は共通次数スカラーに当たる。
従って元の表示はたしかに相対反標準sectionから得られる。
矢は既述のsection再構成で全て基底変更に対応し、
相対Serre自然性により反標準作用との可換性が従う。
共通次数スカラーを除いて $\mathrm{PGL}_4$ 商を得る。
\end{proof}

\subsection{二次包絡の三次平坦性判定}
\begin{lemma}\label{lem:h3-degree-three-flat-koszul}
$A$ を可換局所環、$k$ を剰余体とする。
有限自由加群 $V$ と直和因子 $Q\subset V^{\otimes2}$ から
$B=T_A(V)/(Q)$ を作り、$B_k$ がKoszulであると仮定する。
このとき $B_3$ が $A$ 上自由であることと、
$B_n$ が全ての $n$ で $A$ 上自由であることは同値である。
$A$ のnoetherian性や完備性は仮定しない。
\end{lemma}
\begin{proof}
$B_3$ が自由とする。挿入写像
\[
 Q\otimes_A V\oplus V\otimes_A Q\longrightarrow V^{\otimes3},
 \qquad (a,b)\longmapsto a-b
\]
の余核は $B_3$ である。
従って像は有限自由の直和因子であり、
核 $W_3=(Q\otimes_A V)\cap(V\otimes_A Q)$ も有限自由で、
その形成は底変更と可換する。

関係の生成次数から、任意の $n\ge3$ について次の列がある。
\[
 W_3\otimes_A B_{n-3}
 \xrightarrow{d_3} Q\otimes_A B_{n-2}
 \xrightarrow{d_2} V\otimes_A B_{n-1}
 \longrightarrow B_n\longrightarrow0.
\]
各微分は左因子の最後の一次元を右側の積へ移す。
$W_3\subset V\otimes_A Q$ から $d_2d_3=0$ が従う。
末尾の右完全性は、次数 $n$ の関係を
$Q\otimes_A V^{\otimes(n-2)}$ と
$V$ に次数 $n-1$ の全関係を掛けた部分との和で書くことによる。
剰余体上ではこの列は右Koszul解消の一部であり、
$Q_k\otimes_k B_{k,n-2}$ で完全である。

$n$ 未満の各次数で自由性を仮定する。
$\ker(d_2\otimes_A k)$ の基底を
$d_3\otimes_A k$ の像から選び、$d_3$ の像へ持ち上げる。
これらを $Q\otimes_A B_{n-2}$ の基底へ補うと、
$d_2$ は前半で恒等的に零であり、
後半の像は剰余体上で線型独立だから
$V\otimes_A B_{n-1}$ の基底へ補える。
従って余核 $B_n$ は自由である。
$n\le3$ の仮定から帰納法が成立する。
逆向きは $n=3$ を取ればよい。
\end{proof}

\begin{theorem}\label{thm:h3-finite-degree-normal-incidence}
四変数の幾何学的H3表示を含む正規三次incidenceの形式近傍は、
次の有限次数条件だけで構成できる。
\begin{enumerate}
 \item $Q\subset V^{\otimes2}$ を階数6の部分束とし、
 $B_3,B_4$ の階数を $20,35$ に固定する。
 \item 二次双対の次数3,4の階数を $4,1$ に固定し、
 次数5を零とする。topへの積対合を完全とする。
 \item $B_3$ の切断直線 $\mathcal G$ に対し、$B_4$ において
 $Gx=\mu_B^{-1}(x)G$ を全ての $x\in V$ について課す。
\end{enumerate}
ここで階数条件は余核の局所自由性を表す階数stratumで課す。
幾何学的H3点でのこのincidenceの形式完備化を
$\mathrm{PGL}_4$ で割ると、
系\ref{cor:h3-artin-intrinsic-formal-atlas}の内在変形群oidを得る。

さらに任意の底からこのincidenceへの射で、
各幾何学的fiberがH3表示であるものを取る。
このとき $B$ と $R=B/(\mathcal G)$ は、いずれも各次数で有限局所自由である。
正則正規三次の列と全次数の形成は任意底変更と可換する。
\end{theorem}
\begin{proof}
有限階数stratumは各関係挿入行列で所定の階数より大きい
小行列式を零にし、所定階数の小行列式の非消滅開を取って構成する。
$B_3$ の場合は $64\times48$ 行列の階数44である。
二次双対の次数1,2は最初から階数4,6であり、
指定した残りの階数と積対合で相対Frobenius中山作用が定義される。
切断直線は $\mathbb P(B_3)$ の普遍直線に対応する。
第三条件は $\mathcal G\otimes_A V\to B_4$ の二写像の差が零という
有限個の閉方程式である。

任意の底の点の局所環において、幾何学的Koszul性は剰余体へ忠実平坦降下する。
補題\ref{lem:h3-degree-three-flat-koszul}により各 $B_n$ は局所環上自由となる。
第三条件で $\mathcal G$ は正規である。
左右の倍写像は幾何学的fiber上で単射であり、
各次数で有限自由加群間の分裂単射となる。
よってその余核 $R_n$ も自由である。
$B_n,R_n$ は有限提示なので、これらの局所環での自由性から有限局所自由性を得る。
倍写像の列は局所分裂し、商の形成は任意底変更と可換する。

Artin局所底では剰余体の一点がH3であれば前の議論が適用できる。
逆に全次数平坦な内在Artin変形は既に示したNakayama条件を満たし、
相対Koszul双対の必要な各次数は自由である。
従ってこの有限次数incidenceが、従来の全次数上界による
形式完備化と同じ群oidを表す。
\end{proof}

従って内在的Artin変形で二次包絡の全次数平坦性を調べるには、
反標準環の二次関係を挿入する一つの三次行列を調べれば十分である。
幾何学的H3点の集合がこの有限次数incidence上で通常の開になるかは、
別の相対許容性の問題として残る。

\subsection{一般局所底での平坦対象とtrait上の大域次元}
\begin{proposition}\label{prop:h3-noetherian-local-lift}
可換noetherian局所環 $A$ 上の連結次数付き環 $T$ は
各次数で有限自由であり、剰余体上の環 $T_k$ は左右noetherianであるとする。
このとき $T$ は左右noetherianである。
\end{proposition}
\begin{proof}
まず $A$ を完備とする。極大イデアル $\mathfrak m$ による随伴環は
$T_k\otimes_k\operatorname{gr}_{\mathfrak m}A$ である。
第二因子は $k$ 上有限生成な可換代数なので、随伴環は左右noetherian。
斉次右イデアル $J$ については、交叉で濾過した随伴イデアルの
有限個の二重斉次生成元を $J$ へ持ち上げる。
任意の斉次元について、初期式を順に消すと、
補正係数は各固定次数の自由 $A$ 加群の中で $\mathfrak m$ 進収束する。
元の斉次元は持ち上げた有限個の元と収束した係数の積の和になる。
分離性は各次数の有限自由性から従う。
従って $T$ は斉次右noetherian。
次に通常の次数フィルターを任意の右イデアルに使い、
最高次数を有限回下げることにより右noetherian性を得る。
左も同じである。
一般の $A$ では $\widehat A$ による底変更後にこの結果が適用できる。
任意の右イデアルの昇鎖をそこへ忠実平坦に移し、
停止してから収縮すれば元の鎖も停止する。
\end{proof}

\begin{theorem}\label{thm:h3-relative-flat-coherent-serre-local}
有限次数incidence
（定理\ref{thm:h3-finite-degree-normal-incidence}）から
noetherian局所底 $A$ 上の表示を取り、幾何学的剰余表示がH3であるとする。
$\mathcal F,\mathcal H$ がその $\qgr R$ の底上平坦なcoherent対象なら、
$\mathbf R\operatorname{Hom}(\mathcal F,\mathcal H)$ は $A$ 上完全で
コホモロジーが $0,1,2$ に限られ、剰余体および任意の平坦なnoetherian底変更と可換する。
全てのこのような対象は $\pi R(i)$ で生成される厚い部分圏に入る。
相対Serreはこの厚い圏上で
\[
 S_A=(-1)\otimes_A\mathcal L_{\rm AS}[2]
\]
である。
\end{theorem}
\begin{proof}
三次平坦性判定と正則切断により $B,R$ は次数別に有限自由で、
命題\ref{prop:h3-noetherian-local-lift}から左右noetherianである。
相対Koszul解消とそのFrobenius対合も成立する。
幾何学的剰余体での主張は忠実平坦な体拡大から降ろせる。

まず $\mathcal H$ を次数別に有限自由な有限生成加群 $N$ で表せる。
任意のcoherent代表を取り、その $\operatorname{Tor}_1^A(-,k)$ を考える。
有限生成項を持つ $k$ の $A$ 自由解消の最初の三項から、
これは有限生成次数付き加群である。
torsion商はexactで $A$ 加群とのtensorと可換するので、
$\mathcal H$ の平坦性からこのTorの像は零。
従って十分高い各次数でTorが零となり、有限 $A$ 加群の局所平坦性判定から
それらの次数では代表が自由となる。高い次数で切り詰めれば $N$ が得られる。
$\mathcal F$ にも同じ議論が適用できる。

$N$ を $B$ 加群と見る。剰余体上の有限生成次数付き自由解消を、
各段の斉次生成元を持ち上げることで $A$ 上へ持ち上げる。
剰余体上の全射は次数別のNakayama補題で全射となり、
核も次数別に自由なので底変更で解消性が保たれる。
剰余体上の包絡は大域次元4なので、最後の自由核もその基底を持ち上げればよい。
従って $N$ は $B$ 上長さ4以下の有限自由shift解消を持つ。
$B$ の相対局所双対から、$\pi B(i)$ の各次数の導来sectionは
有限な $A$ 完全複体であり、底変更と可換する。
$R$ 加群 $N$ の局所torsionを求めるときは、
$B/B_{\ge n}$ の商への変換の第一次Torが三次以上離れた遷移で零になるため、
$R$ と $B$ の飽和sectionは同じである。
これは定理\ref{thm:h3-artin-relative-serre-exists}の有限局所cohomology次元の議論である。
従って $N$ の自由解消から $\mathbf R\omega\pi N$ の各次数が底上完全である。
局所torsionは $B/B_{\ge n}$ の長さ4の自由解消で計算されるので、
この比較は剰余体および平坦な底変更と可換する。

$\mathcal F$ の自由shift解消は $R$ 上で取り、各段に有限生成自由加群を使う。
次数別に分裂完全であるため、剰余体や平坦な底変更後も解消である。
$\mathcal H$ への導来Homを取ると、各項は上述の $N$ の飽和sectionで、
底上のコホモロジー振幅は一様に $0$ から $3$ までにある。
従って全導来Homは下に有界で、各有限コホモロジー範囲において
有限生成射影複体で計算され、底変更と可換する。
各自由shiftからの導来Homは、有限自由の飽和複体で表せる。
剰余体上の振幅が $[0,3]$ にあるので、局所的な単元対の消去により
この範囲の自由複体を選べる。
解消の有限切り詰めを一項ずつ伸ばして導来Homを取ると、
第 $j$ 自由項の寄与はその飽和複体を $[-j]$ だけ移したものになる。
各三角を有限自由複体間のchain mapの錐で表すと、
切り詰め間の射を射影で表す逆系が得られる。
各固定次数には有限個の項のみしか寄与せず、
その逆極限で全導来Homを下に有界な有限自由項の複体として表せる。
この複体の単元対は下の次数から順に分離できる。
剰余体上では極小複体の微分は零であり、
幾何学的剰余圏は大域次元2なので、その自由項は $0,1,2$ に限られる。
従って全導来Homはその範囲の完全複体である。

正規三次自己交叉の変換射の二乗は、
底上平坦な $\mathcal F$ とその捩れshift間の第四Extに入るため零である。
包絡への押出しは上記の有限解消を持つから、
二つの錐と直和因子の議論により $\mathcal F$ は $\pi R(i)$ の厚い圏に入る。
その上の相対Serreは、局所化複体の双対
$D_A D_R\simeq D_R(-1)\otimes_A\mathcal L_{\rm AS}[2]$
から得られる。定理\ref{thm:h3-artin-relative-serre-exists}のこの計算に使うのは、
左右の有限局所コホモロジー次元、相対Koszul対合、内部次数別有限自由性、
および飽和の冪等性である。これらは今の底でも成立する。
中山twistが消えるのはincidenceの正規化線型式による。
\end{proof}

\begin{corollary}\label{cor:h3-trait-homological-generization}
前定理の底が離散付値環なら、幾何学的一般fiberでも
$B,R$ は左右noetherianな整域、$B$ はKoszul AS4であり、
$\qgr R$ は大域次元2で $S=(-1)[2]$ を持つ。
構造対象の例外性、反標準section、Hilbert値、一次生成と豊富性も保たれる。
さらに系\ref{cor:h3-rank-from-classified-smooth-domain}の分類された整域のrank判定により、
幾何学的一般fiberもH3である。
\end{corollary}
\begin{proof}
付値の随伴環は剰余環に一つの中央変数を加えたものなので、
剰余表示の整域性から $B,R$ の整域性が従う。
Koszul解消、Frobenius対合および正則正常三次は底変更する。
noetherian性は局所化で保たれる。
一般fiberの有限生成次数付き加群は、その斉次生成元を選び、
元の底の環で張る格子によって延長できる。
この格子はねじれがないので離散付値環上平坦である。
任意の二対象をこうして延長して前定理を適用すれば、一般fiberで
三次以上のExtが零となる。剰余環を幾何化する議論と同じく、
一般fiberも任意の有限拡大を経て代数閉包へ移せる。
実際、有限表示の加群とその有限図式は有限体拡大上に定義され、
その拡大に付値を延長した離散付値環へ同じ証明を適用できる。
代数閉包上のnoetherian性は次の格子判定で従う。
延長付値環を $V'$ とし、値群が離散であることは仮定しない。
代数閉包上の斉次右イデアル $I$ の各次数を、$V'$ 上の自由格子と交叉させる。
有限次元部分空間と標準格子との交叉は有限自由な直和因子である。
実際、一つの非零ベクトルの座標のうち値が最小のものを選び、
その座標を $1$ に規格化する。その座標による射影で残りを帰納すればよい。
従って交叉の剰余は剰余環の斉次右イデアルを与える。
幾何剰余環で有限個の斉次生成元を選び、元の交叉へ持ち上げる。
各次数の交叉は有限自由なので、次数別Nakayama補題により
持ち上げた生成元が交叉を生成し、従って $I$ を生成する。
最高次数による帰納法で任意の右イデアルについてもnoetherianとなる。
左も同じである。
この上で、有限体拡大上のExt消滅は体の平坦底変更で代数閉包へ移る。

大域次元の下界2は
$\operatorname{Ext}^2(\pi R,\pi R(-1))\simeq k$ から従う。
体上の相対Serreの式および飽和sectionの計算を適用し、
例外性とsectionの各主張を得る。
有限生成次数付き加群の包絡上の有限解消から $\chi$ が従い、
Artin--Zhang豊富性を得る。
\end{proof}

\begin{remark}
任意底の表示の次数スカラーを消す際、
命題\ref{prop:h3-presentation-canonical-degree-descent}の重み一の線には
$\mathcal L_{\rm AS}=W_4\otimes\mathcal G^\vee$ も使える。
$W_4$ は重み4、$\mathcal G$ は重み3なので同じ降下が働く。
この選択による商環の次数 $n$ は $R_n\otimes\mathcal L_{\rm AS}^{-n}$ であり、
相対Serreから計算した反標準sectionそのものに対応する。
従ってこの標準化では相対orientationの自明化を別の標識にする必要がない。
\end{remark}
